\chapter{Introduction}

\section{Background and motivation}

Academic and domain-specific terminology often admits multiple valid lexical formulations for the same underlying concept, and semantic relatedness does not necessarily imply strict synonymy: a term can be topically or associatively close to another without being an acceptable substitute for it. Retrieving the specific synonym intended for a given term from a controlled terminology inventory is therefore a non-trivial ranking problem, distinct from general semantic-similarity search.

This thesis studies that problem empirically, comparing lexical retrieval, dense embedding retrieval, Large Language Model (LLM)-based query expansion, and LLM-based reranking as different mechanisms for identifying the intended synonym of an English academic or domain term, and comparing a fixed-role multi-agent pipeline against a more autonomous, dynamic tool-calling alternative for orchestrating these mechanisms.

\section{Research problem and task formulation}

Given an input academic or domain term, the system ranks candidates from a fixed terminology inventory and is evaluated according to the position of the row-specific annotated synonym. The task is formulated as \textbf{transductive, closed-pool synonym ranking}: a global candidate vocabulary is available in full at evaluation time, including candidate strings associated with validation and test rows, while the correct mapping between an individual input term and a vocabulary entry remains hidden from the system until after a prediction has been produced. This formulation differs from unrestricted synonym generation, since every returned answer must belong to the fixed inventory; the reported results measure ranking and selection quality within that known inventory and should not be interpreted as evidence of the ability to generate synonyms outside it.

The candidate vocabulary is derived from the dataset prepared for the task. One degenerate identity pair -- an input term recorded as its own synonym -- is then removed to form the final evaluation dataset (Chapter~\ref{chap:data}), which contains 969 rows. The candidate vocabulary itself has a verified cardinality of 921 candidate synonyms.

\section{Research objectives}

The project was structured around the following objectives, each evaluated empirically under a common protocol:

\begin{itemize}
  \item evaluate lexical retrieval and dense embedding retrieval as baseline approaches to closed-pool synonym ranking;
  \item evaluate LLM-based query expansion, and deterministic weighted multi-query fusion as a mechanism for controlling its contribution;
  \item evaluate LLM-based reranking of a retrieved shortlist, including its sensitivity to reranker model capability;
  \item compare a fixed-role, centrally orchestrated multi-agent architecture against a dynamic, single-agent tool-calling architecture under matched downstream retrieval and reranking components;
  \item examine the accuracy--latency--cost trade-offs across the resulting configurations.
\end{itemize}

\section{Proposed approach}

The system evaluated as the primary configuration is a \textbf{fixed-role, centrally orchestrated multi-agent pipeline}. A Generator Agent, an LLM-based component, produces query variants and paraphrases of the input term. Deterministic retrieval and weighted fusion -- supporting software components, not agents -- use these variants to score and shortlist candidates from the closed vocabulary using dense embeddings. A Verifier/Reranker Agent, a second LLM-based component, then evaluates and reorders that shortlist to produce the final prediction. Neither retrieval nor fusion negotiates, plans, or allocates tasks; coordination between the two agents is centralized rather than decentralized.

As an alternative, the project also implements and evaluates a \textbf{dynamic tool-calling architecture}, in which a single controller agent decides at runtime, rather than through a fixed expansion strategy, which retrieval queries to issue and how many retrieval calls to perform before finalizing its candidate list. Retrieval remains a tool invoked by this controller, not an agent in its own right. The two architectures are compared experimentally under a matched downstream configuration, sharing the same retrieval model and final reranker (Chapter~\ref{chap:results}).

\section{Thesis structure}

The remainder of this thesis is organized as follows. Chapter~\ref{chap:data} (Data) describes the source dataset and its documented provenance, the construction of the project-specific dataset and closed candidate vocabulary, preprocessing and normalization, the train/validation/test partitioning protocol, and the resulting task and leakage boundaries. Chapter~\ref{chap:models} (Models) presents the baseline model families, the fixed-role multi-agent pipeline and the dynamic tool-calling architecture in detail, the prompts and interface contracts used by each agent, engineering trade-offs across configurations, and the formal evaluation protocol. Chapter~\ref{chap:results} (Results) reports the main quantitative results, validation ablations, the fixed-versus-dynamic architecture comparison, the accuracy--latency trade-off, error analysis, and statistical interpretation of the reported differences. Chapter~\ref{chap:conclusions} (Conclusions) revisits the research objectives in light of the experimental evidence, synthesizes the methodological choices and results, discusses limitations and engineering recommendations, and proposes future work.
